
\setcounter{footnote}{39}

\SGARefTarget{sga3:iv:subsection:4:5}\subsection*{4.5. The Case of a Topology Coarser than the Canonical Topology}
\addcontentsline{toc}{subsection}{\SGARefLink{sga3:iv:subsection:4:5}{4.5}. A Topology Coarser than the Canonical Topology}

By \SGARefLink{sga3:iv:definition:4:3:6}{4.3.6} and \SGARefLink{sga3:iv:definition:4:3:8}{4.3.8}, the following conditions on a topology \(\mathcal T\) on
\(\mathscr C\) are equivalent:
\begin{enumerate}
  \item[(i)]\SGARefTarget{sga3:iv:subsection:4:5:item:i:01} \(\mathcal T\) is coarser than the canonical topology of
    \(\mathscr C\);
  \item[(ii)]\SGARefTarget{sga3:iv:subsection:4:5:item:ii:01} every representable presheaf is a sheaf for \(\mathcal T\);
  \item[(iii)]\SGARefTarget{sga3:iv:subsection:4:5:item:iii:01} every covering sieve for \(\mathcal T\) is universal effective
    epimorphic.
\end{enumerate}
If \(\mathcal T\) is defined by a pretopology
\(S\mapsto\mathscr R(S)\), these conditions are also equivalent to:
\begin{enumerate}
  \item[(iv)]\SGARefTarget{sga3:iv:subsection:4:5:item:iv:01} every family belonging to \(\mathscr R(S)\) is universal
    effective epimorphic.
\end{enumerate}

When these conditions hold, the canonical functor
\(\mathscr C\to\widehat{\mathscr C}\) factors through a functor
\[
  j_{\mathscr C}=j:\mathscr C\longrightarrow\widetilde{\mathscr C};
\]
we also write \(j(S)=\widetilde S\).\footnote{Editor's note: we shall also
write \(\mathbf h_S=\widehat S\); cf. the first commutative diagram in
\SGARefLink{sga3:iv:corollary:4:5:4}{4.5.4}.}

\medskip
\noindent\SGARefTarget{sga3:iv:proposition:4:5:1}\textbf{Proposition 4.5.1.}
\textit{The functor \(j:\mathscr C\to\widetilde{\mathscr C}\) is fully
faithful and commutes with arbitrary inverse limits. In particular, it is
left exact and therefore preserves algebraic structures defined by finite
inverse limits.}

This follows immediately from the commutative factorization
\[
\begin{array}{ccc}
 &\mathscr C&\\[-0.5ex]
 {\scriptstyle j}\swarrow&&\searrow\\[-0.5ex]
 \widetilde{\mathscr C}&\xrightarrow{\ i\ }&\widehat{\mathscr C}
\end{array}
\]
and from \SGARefLink{sga3:iv:theorem:4:4:1:item:i:01}{4.4.1(i)}.

Before giving further properties of \(j\), we must define the topology
induced on a category \(\mathscr C/S\). We no longer suppose here that the
given topology is necessarily coarser than the canonical topology. A sieve
of an object \(T\) in \(\mathscr C\), together with a morphism \(T\to S\),
naturally defines a sieve of \(T\) in \(\mathscr C/S\), because \SGARefLink{sga3:iv:definition:4:1:2}{the
definition} of a sieve of \(T\) depends only on the category
\[
  \mathscr C/T=(\mathscr C/S)/T.
\]
For example, if the sieve is defined by a family \(\{T_i\to T\}\), its image
in \(\mathscr C/S\) is defined by the same family, now regarded as a family
of morphisms in \(\mathscr C/S\).

Declaring these images of covering sieves to be covering defines a topology
on \(\mathscr C/S\), called the \emph{topology induced} by the given
topology. In the terminology of \SGARefLink{sga3:iv:bib:as}{[AS]}, \SGARefLink{sga3:iv:proposition:2:3}{2.3}, it is the coarsest topology on
\(\mathscr C/S\) for which the canonical functor
\[
  i_S:\mathscr C/S\longrightarrow\mathscr C
\]
is a comorphism.\footnote{Editor's note: that is, for every object
\(T\to S\) of \(\mathscr C/S\), every covering sieve \(R'\) of
\(i_S(T\to S)=T\), when regarded as a sieve of the object
\((T\to S)\in\operatorname{Ob}(\mathscr C/S)\), is covering.}
By definition, the identifications
\[
  (\mathscr C/S)/T=\mathscr C/T
\]
respect these topologies.
